% ============================================================
%  Figure : Hiérarchie des classes de modèles
%  Compiler : pdflatex fig_hierarchy.tex
% ============================================================
\documentclass[tikz, border=8pt]{standalone}

\usepackage[utf8]{inputenc}
\usepackage[T1]{fontenc}
\usepackage{xcolor}
\usepackage{tikz}
\usetikzlibrary{positioning, arrows.meta}

\definecolor{pythonblue}{HTML}{356D9C}
\definecolor{pythonyellow}{HTML}{FED03E}
\definecolor{pythondark}{HTML}{1E4068}
\definecolor{pythonlight}{HTML}{D6E8F5}

% Macro pour dessiner une boîte de classe UML à 3 compartiments
% #1 = nom du noeud, #2 = position, #3 = classe, #4 = attributs, #5 = méthodes
\newcommand{\umlclass}[5]{
  \node (#1) at #2 {%
    \renewcommand{\arraystretch}{1.3}
    \begin{tabular}{|p{5.2cm}|}
      \hline
      \cellcolor{pythonlight}\centering\small\ttfamily\bfseries #3 \\
      \hline
      \small\ttfamily #4 \\
      \hline
      \small\ttfamily #5 \\
      \hline
    \end{tabular}
  };
}

\newcommand{\umlabstract}[5]{
  \node (#1) at #2 {%
    \renewcommand{\arraystretch}{1.3}
    \begin{tabular}{|p{5.2cm}|}
      \hline
      \cellcolor{pythonlight!60}\centering\small\ttfamily\itshape\bfseries
        $\langle\langle$abstract$\rangle\rangle$ \\ [-4pt]
      \cellcolor{pythonlight!60}\centering\small\ttfamily\itshape #3 \\
      \hline
      \small\ttfamily\itshape #4 \\
      \hline
      \small\ttfamily\itshape #5 \\
      \hline
    \end{tabular}
  };
}

\begin{document}
\begin{tikzpicture}[
  inherit/.style={-{Triangle[open, length=9pt, width=7pt]}, thick, pythonblue},
  assoc/.style={-{Stealth[length=7pt]}, thick, pythonblue!70, dashed}
]

% ---- BaseRegressor (abstract) ----
\umlabstract{br}{(0, 0)}{BaseRegressor}{%
  + config: Config \\
  + model: sklearn \\
  + scaler: RobustScaler%
}{%
  + fit(X, y) \\
  + predict(X) \\
  + evaluate(X, y) \\
  + compute\_r2\_score(...) \\
  + compute\_rmse(...) \\
  + compute\_mae(...) \\
  + check\_residuals\_normality(...) \\
  -- \textit{\_fit(X, y)  [abstract]}%
}

% ---- LinearRegressor ----
\umlclass{lr}{(-4, -7)}{LinearRegressor}{%
  (hérite de BaseRegressor)%
}{%
  + get\_coefficients() \\
  -- \_fit(X, y)%
}

% ---- EnsembleRegressor ----
\umlclass{er}{(4, -7)}{EnsembleRegressor}{%
  + model\_type: str \\
  \quad ('random\_forest' | 'gradient\_boosting')%
}{%
  + get\_feature\_importance(...) \\
  -- \_fit(X, y)%
}

% ---- Config ----
\node (cfg) at (8.5, 0) {%
  \renewcommand{\arraystretch}{1.3}
  \begin{tabular}{|p{4.5cm}|}
    \hline
    \cellcolor{pythonyellow!40}\centering\small\ttfamily\bfseries Config \\
    \hline
    \small\ttfamily
    + RANDOM\_STATE: int \\
    + TEST\_SPLIT\_RATIO: float \\
    + RF\_N\_ESTIMATORS: int \\
    + RF\_MAX\_DEPTH: int \\
    + GB\_LEARNING\_RATE: float \\
    + ... \\
    \hline
    \small\ttfamily
    + \_validate() \\
    + save\_to\_file(path) \\
    \hline
  \end{tabular}
};

% Héritage
\draw[inherit] (lr.north) -- ++(0, 1.0) -| (br.south);
\draw[inherit] (er.north) -- ++(0, 1.0) -| (br.south);

% Association Config
\draw[assoc] (cfg.west) -- (br.east)
  node[midway, above, font=\tiny\itshape, color=pythonblue]{injecté};

\end{tikzpicture}
\end{document}
